\section{跨模块工业级技术评价基线}

本章规定本手册所述模块进入工程算例、批量研究或生产服务前必须共同满足的评价基线。
具体模型公式、算法步骤、选项和终止判据以正文相应章节为准；本章不扩大源码已经实现的范围。

\subsection{输入、输出、边界条件与数据结构审计}
输入审计应逐项检查问题维度一致性、系数有限性、上下界合法性（含 $\ell\le u$）、
变量类型标记（连续/整数/二进制）、锥成员归属与选项枚举值；非法输入必须在进入求解内核前由
\srcpath{src/engine/util/problem\_validation.cpp} 一类入口拒绝并返回结构化错误。
输出审计必须记录求解状态（最优/不可行/无界/时限/数值异常且相互区分）、目标值、
原始与对偶残差、间隙、后端名称、容差、迭代/节点计数和告警。
空问题、零行零列、退化约束、固定变量、重复索引与不可用可选后端均属于边界测试集。

\subsection{工程假设与数值稳定性分析}
模型使用的凸性、光滑性、约束规范、中心路径存在性、基非奇异性等假设必须与所求解的问题类匹配。
数值评价至少包含数据缩放（scaling）、原始/对偶不可行度、互补间隙、KKT 残差、
基或因子分解的条件数/枢轴质量、步长/阻尼、迭代停滞与容差复现性。
若采用正则化、扰动、惯性校正、回退后端或近似（如容差放宽），应同时报告触发条件、参数值、
对主指标的敏感性以及是否改变结果口径；不得用``已返回结果''替代最优性或收敛性证明。

\subsection{复杂度与资源分析}
令 $m$、$n$、$n_z$ 分别表示约束数、变量数与矩阵非零元数。
报告应分别给出建模、预求解、求解、后处理的时间与内存。
单纯形单次迭代的稠密代价为 $O(m^2)$ 量级，实际取决于基因子更新与定价策略；
内点法每步代价由 KKT 稀疏因子分解主导，应报告结构非零数、填充率和因子分解峰值内存；
MILP 不给出虚假的多项式最坏界，而报告节点数、间隙、时限和实测扩展曲线。
复杂度结论必须基于固定硬件、固定后端、固定构建配置和可复现算例族。

\subsection{异常工况处理}
不可行、无界、数值失败、时限、内存不足和后端缺失必须相互区分并沿 API 如实上报；
不得把数值失败静默改写为不可行，不得把次优解标为全局最优，不得用零值或 NaN 伪装未知量。
异常路径必须保证问题对象可恢复、可重试；回退（fallback）到替代后端时必须在结果或日志中
声明实际使用的求解器。

\subsection{与其他模块的接口关系}
接口链路遵循``AML/文件输入 $\to$ 问题校验 $\to$ 预求解 $\to$ 求解内核或外部后端
$\to$ 后处理/解回写''。跨模块传递时保留变量/约束的稳定身份、索引空间、符号约定与单位；
消费 SCUC 等上层应用结果前，先检查其求解状态与容差口径。Python 绑定层只做语义保持适配，
不重新解释求解结果。

\subsection{测试与验证建议}
最低验证矩阵包括：解析可解小例、标准算例集（如 NETLIB、MIPLIB 2017）、边界/异常输入、
算法或后端交叉校核、序列化往返、重复运行确定性和规模扩展测试。
数值算法还应加入残差独立重算；近似路径应与高保真路径比较并给出误差分布，而非只报告平均值。
回归报告必须列明提交哈希、构建配置、算例版本、容差、通过/失败/跳过数及未验证范围。

\subsection{工业级评估指标}
\begin{itemize}
  \item \textbf{正确性}：KKT 残差、约束最大违反、互补间隙、解析/基准目标误差、不可行证书正确性；
  \item \textbf{鲁棒性}：标准算例集收敛率、不可行/无界识别率、数值失败率、容差敏感性；
  \item \textbf{性能}：端到端时延、迭代/节点计数、峰值内存、并行效率与规模增长斜率；
  \item \textbf{可追溯性}：选项快照、后端与版本记录、容差口径、告警可定位率；
  \item \textbf{工程有效性}：与参考求解器或历史基线的目标值/时间偏差，及其对业务指标的影响。
\end{itemize}
指标应预先给出验收阈值和失败处置；未达到阈值时只能标记为实验性或受限适用。

\subsection{适用范围、局限性与后续改进}
本手册只评价当前源码覆盖的问题类、后端和选项，不对未实现的问题类、未安装求解器或未执行的
外部验证作保证。后续改进应按``缺口 $\to$ 可量化指标 $\to$ 源码/测试变更 $\to$ 独立验证''闭环，
优先消除会造成错误结论的语义缺口，其次提升数值鲁棒性与性能，最后扩展问题类范围。
